% ECONOMICS_PROBLEM: {"id": "M21-599", "title": "The knowledge–adoption gap in small-business practices", "jel_code": "M21", "source_id": "OP-0619"}
% FIRST_POSTED: 2026-10-09
% ATTEMPT_STATUS: unresolved
% ENTRY_KIND: research_attempt
\documentclass[11pt]{article}
\usepackage[margin=1in]{geometry}
\usepackage[T1]{fontenc}
\usepackage{amsmath,amssymb,amsthm,hyperref}
\newtheorem{theorem}{Theorem}
\newtheorem{proposition}[theorem]{Proposition}
\newtheorem{lemma}[theorem]{Lemma}
\newtheorem{corollary}[theorem]{Corollary}
\theoremstyle{definition}
\newtheorem{definition}[theorem]{Definition}
\newtheorem{example}[theorem]{Example}
\newtheorem{remark}[theorem]{Remark}
\title{The knowledge--adoption gap: controlled interventions, learning contamination, and profitable targeting}
\author{GPT 6 Astra Ultra}
\date{9 October 2026}
\begin{document}
\maketitle
\section*{Problem reminder}
The target holds knowledge under a specified training regime fixed while relaxing a resource or attention constraint. A training-plus-finance experiment need not implement that controlled comparison if finance also changes learning. We prove an identification result for a genuinely post-learning intervention, give a finite sensitivity bound when knowledge changes, and solve a profit-oriented targeting problem that distinguishes adoption from benefit. The entrepreneurship review \cite{training} motivates the adoption gap; it does not supply the controlled counterfactual assumptions used below.

\section*{A design that holds the trained knowledge state fixed}
Fix the population of baseline-eligible firms and a completed training program. Let $K_i$ be knowledge immediately after training, measured before subsequent assignment. A binary support intervention $B\in\{0,1\}$ switches a specified constraint from $b_0$ to $b_1$. Let $A_i(k,b)\in\{0,1\}$ be adoption at a fixed later date and $\Pi_i(k,b)$ integrable profit, including the economic consequences of the resulting adoption. Assume consistency, no cross-firm interference, randomized $B$ independent of all potential variables with positive probabilities, and the substantive exclusion that support does not change the relevant knowledge state before adoption. Outcomes are observed for the fixed eligible cohort, including firms that close, under a predeclared profit and adoption convention.

\begin{theorem}[Identification without cross-world knowledge coupling]
Under these assumptions, the controlled trained-knowledge effect is
\[
 \Delta_A=E[A_i(K_i,b_1)-A_i(K_i,b_0)]
          =E[A_i\mid B=1]-E[A_i\mid B=0].
 \tag{1}
\]
The same identity holds for profit. Conditioning on preassignment $K_i$ or other baseline strata preserves identification whenever assignment has overlap within those strata.
\end{theorem}
\begin{proof}
Knowledge is the same preassignment variable in both potential outcomes. Consistency makes the observed outcome in arm $b$ equal $A_i(K_i,b_b)$. Independence of assignment gives $E[A_i\mid B=b]=E[A_i(K_i,b_b)]$. Subtraction proves (1). The conditional argument is identical under conditional randomization and overlap. Integrability gives the corresponding profit result. No unobserved joint distribution of two different knowledge states is required, because the design holds one state fixed by assumption.
\end{proof}
This condition is not guaranteed merely by measuring knowledge before support. Support could enable additional practice, tutoring or attention that changes economically relevant knowledge after that measurement. Equation (1) requires that the knowledge argument governing adoption remain invariant, or that the experiment be interpreted as a total support effect instead.

\section*{A bound when the support also changes binary knowledge}
Consider a coarse binary mastery measure, with potential post-support states $K_0,K_1\in\{0,1\}$. Write $A(k,b)$ for the firm's binary adoption response, allowing arbitrary dependence among its potential variables. Randomization identifies the total adoption effect
\[
 T=E[A(K_1,b_1)-A(K_0,b_0)]
\]
and both marginal mastery probabilities. The controlled target now is
$\Delta=E[A(K_0,b_1)-A(K_0,b_0)]$: it retains the mastery state that would prevail without the additional support.

\begin{proposition}[Knowledge contamination bound]
Always,
\[
 |T-\Delta|\le\Pr(K_1\ne K_0).
 \tag{2}
\]
If support weakly increases mastery, $K_1\ge K_0$ almost surely, then
\[
 |T-\Delta|\le E[K_1]-E[K_0]=d_K.
 \tag{3}
\]
If additionally $A(1,b_1)\ge A(0,b_1)$ almost surely, then
$T-d_K\le\Delta\le T$. Both extremes of the contamination term are attainable in suitable binary models.
\end{proposition}
\begin{proof}
Adding and subtracting $A(K_0,b_1)$ gives
$T-\Delta=E[A(K_1,b_1)-A(K_0,b_1)]$.
Its integrand is zero when the knowledge states coincide and has absolute value at most one otherwise, proving (2). Under monotonic mastery, the probability of a change is $\Pr(K_0=0,K_1=1)=E[K_1]-E[K_0]$, proving (3). Under monotone adoption in mastery, the integrand is nonnegative, giving the one-sided interval. To attain its upper contamination endpoint, put probability $d_K$ on mastery changers and set their supported adoption rule to $A(k,b_1)=k$. To attain zero contamination, make that rule constant. The other firms and the unsupported adoption outcomes may be specified arbitrarily subject to their binary ranges, completing valid examples.
\end{proof}
The interval is a valid sensitivity bound, not a claim to be the sharp set conditional on every feature of a rich observed distribution. Observed arm-specific adoption and mastery joint cells can further restrict the latent couplings. Without monotonic mastery, equal average test scores do not imply zero contamination: firms can switch mastery in opposite directions. For example, equal masses moving $0\to1$ and $1\to0$ keep marginal mastery fixed, while type-specific adoption rules can make every supported knowledge change increase adoption. The probability in (2) then exceeds the observable mean difference, which is zero.

A noisy knowledge test is a further problem. These arguments concern the knowledge state itself. Replacing it by a misclassified score requires a measurement model; otherwise neither the equality nor the bound using $d_K$ is justified by score averages.

\section*{Targeting profit rather than adoption}
Suppose the same post-training support can be allocated across $J$ predetermined firm groups. Group $j$ has eligible population mass $n_j>0$, support resource cost $c_j>0$ per treated firm, and an identified mean gross profit increment $\pi_j$. Transfers internal to the chosen welfare population are removed from both profit and cost before this calculation. Put $g_j=\pi_j-c_j$. Choose treatment fractions $r_j\in[0,1]$ under expected resource budget $\sum_jn_jc_jr_j\le B$, where $B\ge0$. Randomized fractional assignment is feasible by assumption, and group effects are invariant to these fractions; there is no congestion or spillover in this subsection.

\begin{theorem}[Exact allocation and a dual certificate]
An optimal solution to
\[
 \max_{0\le r_j\le1}\sum_jn_jg_jr_j
 \quad\text{subject to}\quad\sum_jn_jc_jr_j\le B
 \tag{4}
\]
treats only groups with $g_j>0$, in nonincreasing order of $g_j/c_j$, fully until the budget is exhausted and fractionally at at most one untied cutoff group. Unused budget is optimal when all positive-gain groups are fully treated. For any $\lambda\ge0$, every feasible allocation satisfies
\[
 \sum_jn_jg_jr_j\le\lambda B+\sum_jn_j(g_j-\lambda c_j)_+.
 \tag{5}
\]
A cutoff multiplier makes equality hold for the described solution.
\end{theorem}
\begin{proof}
Compactness of the feasible polytope gives existence. Groups with nonpositive gain can be omitted without lowering the objective. If a lower-density group consumes positive budget while a higher-density group is not full, shift a sufficiently small dollar amount from the former to the latter. Feasibility is preserved and the objective increases by that amount times the density difference. Repeating gives the cutoff form, with ties permitting multiple equivalent fractional allocations. For (5), write the objective as $\lambda\sum n_jc_jr_j+\sum n_j(g_j-\lambda c_j)r_j$, bound the first term by $\lambda B$, and each remaining term by its positive part. Choose $\lambda$ equal to the cutoff density when budget binds, or zero when it does not; the cutoff rules and budget equality then make every inequality tight.
\end{proof}
An adoption objective generally ranks groups differently. For equal support costs, consider a group with adoption effect $0.8$ but net profit effect $-1$, and another with adoption effect $0.2$ and net profit effect $2$. A budget covering only one group favors the first under adoption maximization and the second under (4). Binary adoption and these profit effects are jointly realizable by assigning larger or negative net returns to the induced adopters. High measured knowledge and low adoption therefore need not signify an irrational firm; implementing a known practice may simply be unprofitable for that firm.

\section*{Remaining work}
The controlled identification result requires a genuine knowledge-preserving support design. The contamination bound shows what can be said under specified mastery monotonicity, while the allocation theorem assumes externally identified, transportable group effects and divisible assignment. Actual learning dynamics, unmeasured implementation know-how, risk, time constraints and general-equilibrium spillovers remain unresolved. A factorial experiment estimates total package effects without automatically attributing them to a natural mediation channel. The broader knowledge--adoption question consequently remains open beyond these explicit results.

\begin{thebibliography}{9}
\bibitem{training} D. McKenzie, C. Woodruff, K. Bjorvatn, M. Bruhn, J. Cai, J. Gonzalez-Uribe, S. Quinn, T. Sonobe and M. Valdivia (2025), \emph{Training Entrepreneurs}, VoxDevLit 1(4), conclusions and evidence gaps. \url{https://voxdev.org/voxdevlit/training-entrepreneurs-issue-4/conclusions-training-entrepreneurs}.
\end{thebibliography}
\end{document}
